\section{CONCLUSION}

This survey systematically reviewed the technical trends and research
directions in Python type analysis.
Python’s dynamic execution model fundamentally conflicts with the fixed and
closed type structures assumed by traditional static analyses, making it
challenging to build practical analyzers by simply adapting type analysis
techniques from other languages.
Consequently, the Python ecosystem has seen the parallel development of a wide
range of strategies, including static analyzers, runtime monitoring systems,
execution trace-based type inference, symbolic execution, constraint solving,
and more recently, LLM-based approaches.

From this analysis, several conclusions can be drawn:
\begin{itemize}
  \item Static analysis alone cannot fully capture Python’s dynamic features.
    Structural typing, dynamic attributes, reflection, and heterogeneous
    containers violate assumptions of conventional static type theory.

  \item Dynamic and hybrid approaches provide high accuracy but cannot fully
    resolve coverage and performance limitations.

  \item LLM-based approaches offer new possibilities but still fail to
    guarantee soundness.
\end{itemize}

Therefore, a hybrid, multi-layered approach that combines the scalability of
static analysis, the precision of dynamic analysis, and the generalization
capability of AI-based inference is essential in Python.
Overall, Python type analysis remains an evolving domain, situated in the
tension between the language’s dynamic design philosophy and the requirements
of static safety, with each technique playing a complementary role in bridging
this gap.
Future research is expected to formalize the interaction between static,
dynamic, and AI-based techniques, exploring new analysis models that balance
soundness, precision, and scalability.
